% ===== Envoi =====
\newpage
\goldrule
\begin{center}
{\large\bfseries\color{rosedeep} Envoi: The Flower on the Path}
\end{center}
\goldrule

\vspace{0.8em}

\noindent
There is a flower.

\vspace{0.5em}

\noindent
It grows in a crack between paving stones, on a path that has been walked many times. It is not rare. It has no name that anyone bothers to remember. If you asked a jeweller what it was worth, they would not understand the question. It has no economic value, no sign exchange value, no position in any hierarchy of marketed objects. By every measure that the dominant framework of luxury recognises, it is nothing at all.

\vspace{0.5em}

\noindent
And yet.

\vspace{0.5em}

\noindent
A person is walking the path---\emph{this} person, on \emph{this} morning, in the particular quality of light that belongs to this specific moment and no other. They are carrying, as they always carry now, the accumulated history of a shared life: months and years of attentive co-presence, of noticing, of learning what the beloved loves. And because of this history---because the relation has cultivated in them an attention shaped by the beloved's specific way of seeing the world---they notice the flower. Not as a botanist notices a specimen, not as a consumer notices a product, but as a lover notices something the beloved would love.

\vspace{0.5em}

\noindent
\emc{In that moment, the flower becomes the most luxurious object in the world.}

\vspace{0.5em}

\noindent
Not because anything about the flower has changed---it is the same flower in the same crack in the same path---but because the person's spatiotemporal structure has become embedded in it. Their existence in this specific moment, their attention shaped by this specific relation, their care directed toward this specific beloved: all of this is now carried by the flower, as an ineliminable feature of what it is. The flower has undergone the third transfiguration. It has become an existence-attestation. It carries, in its modest and perishable materiality, the trace of a specific existence at a specific moment---a coupling degree that no diamond purchased by proxy could ever match.

\vspace{0.5em}

\noindent
\emc{{\cjk 路边那朵花，那个时刻，那个时空。}}

\vspace{0.5em}

\noindent
The flower on the path, that moment, that spatiotemporal structure. This is what she taught me---not by instruction, but by the quality of her own attention to the world. She, who loves the living world, who finds in a wildflower or a quality of light or a bird's call through an open window something more precious than anything that can be bought: she showed me that the deepest luxury is not in what we own but in what we notice, not in what we can purchase but in what we can attest, not in the economic value of the gift but in the genuine embedding of one existence in the life of another.

\vspace{0.5em}

\noindent
She showed me that a flower, in the right moment, given with genuine presence, is more luxurious than any diamond---and that this is not a sentimental consolation for those who cannot afford diamonds but a precise ontological truth about the nature of luxury in intimate life.

\vspace{0.5em}

\noindent
\emc{Nothing is more luxurious than this.}

\vspace{0.5em}

\noindent
Not the most expensive ring, not the rarest gemstone, not the most prestigious brand. The flower on the path, noticed by a person whose attention has been shaped by love, carried home to a specific beloved in a specific moment that will never come again: this is luxury in its truest form, the luxury that the consumer society cannot see because it has no category for high coupling degree at zero economic cost.

\vspace{0.5em}

\noindent
And it is available to anyone---to anyone who has cultivated the sensibility to notice the particular, the relational history to recognise what is shared, and the presence to attest their existence in the modest objects of an ordinary day. The most luxurious thing in the world grows in a crack between paving stones, free for anyone with the eyes to see it.

\vspace{0.5em}

\noindent
\emc{For her, who has those eyes.}

\vspace{0.5em}

\noindent
For the girl from home, who loves culture and travel and the living world, who taught me what luxury is by the simple and inexhaustible example of her attention to the world. The flower on the path was always there. It took her to teach me how to see it.

\vspace{0.5em}

\noindent
And so this paper, like the flower, is an existence-attestation: the embedding of a specific gratitude, formed across a specific shared life, in the modest materiality of words. It is not the most valuable thing I could give. But it carries, I hope, a coupling degree that no more expensive gift could match.

\vspace{2em}
\goldrule

\vspace{1.5em}
\begin{center}
\begin{minipage}{0.72\linewidth}
\centering\itshape\color{rose}
{\cjk 身无彩凤双飞翼，心有灵犀一点通。}\\[0.5em]
\upshape Without the phoenix wings to fly together,\\[0.2em]
our hearts are joined by the unicorn's single horn of light.\\[0.8em]
{\footnotesize\color{gold} {\cjk 李商隐} · \emph{Untitled}}
\end{minipage}
\end{center}

\vspace{1.5em}
\goldrule
